\section{Notation and Pseudocode}

\subsection{Main Notation}
\begin{table}[h]
\centering
\caption{Main notation used in the paper.}
\label{tab:notation}
\small
\begin{tabular}{ll}
\toprule
Symbol & Meaning \\
\midrule
$\gG(\gV,\gE)$ & Task graph \\
$(j,i)$ & Dependency edge from task $j$ to task $i$ \\
$t_i^{sw}, t_i^{hw}$ & Software and hardware execution costs \\
$A_i, A_{\max}$ & Node area cost and total hardware-area budget \\
$C_{ji}$ & Communication cost on edge $(j,i)$ \\
$\vp, \tilde{\vp}$ & Discrete and relaxed placements \\
$\pi^{sw}, \tilde{\vpi}^{sw}$ & Discrete and relaxed software orders \\
$\vs, \vf$ & Discrete start and finish times \\
$\tilde{\vs}, \tilde{\vf}$ & Relaxed start and finish times \\
$M_{\mathrm{DAG}}$ & DAG makespan baseline \\
$M_{\mathrm{static}}$ & Static-priority makespan \\
$M_{\mathrm{learned}}$ & Learned-priority makespan \\
$\widetilde{M}(\tilde{\vp},\tilde{\vpi}^{sw})$ & Relaxed makespan surrogate \\
$\SMax_{\beta}(\cdot)$ & Smooth maximum operator \\
\bottomrule
\end{tabular}
\end{table}

\subsection{Order-Aware Training and Decode}
\paragraph{Algorithm 1: Training of \diffgnnorder.}
\begin{enumerate}
\item Construct node and edge features for the input task graph $\gG=(\gV,\gE)$.
\item Run the shared encoder and predict the relaxed placement vector $\tilde{\vp}$ and software-priority vector $\tilde{\vpi}^{sw}$.
\item Build the soft software permutation $\mP^{sw}$ with Sinkhorn normalization and convert it to the pairwise-before matrix $\mB^{sw}$.
\item Form the effective software-order matrix $\mR^{sw}$ using Equation~\eqref{eq:effective_sw_order}.
\item Evaluate the relaxed start and finish times with the differentiable scheduler and compute $\widetilde{M}(\tilde{\vp},\tilde{\vpi}^{sw})$.
\item Add the area penalty and auxiliary regularizers to obtain $\mathcal{L}_{\mathrm{ord}}$.
\item Update the model parameters by backpropagation.
\end{enumerate}

\paragraph{Algorithm 2: Discrete inference and reporting.}
\begin{enumerate}
\item Decode one or more candidate binary partitions from the relaxed placement scores.
\item Repair area violations and optionally refill unused hardware area.
\item Apply hybrid LSSP local search to the repaired partition.
\item Compute the DAG baseline $M_{\mathrm{DAG}}$ and the static-priority makespan $M_{\mathrm{static}}$.
\item If learned software-priority scores are available, compute $M_{\mathrm{learned}}$ and report $M_{\mathrm{report}}=\min(M_{\mathrm{static}},M_{\mathrm{learned}})$; otherwise report $M_{\mathrm{report}}=M_{\mathrm{static}}$.
\end{enumerate}

\subsection{Synthetic Scaling Generator}
The new synthetic benchmark generator is implemented in \texttt{figures/synthetic.py} and exposed through \texttt{tools/generate\_squeezenet\_like\_synthetic\_graph.py}. It uses the SqueezeNet-TOSA topology as a template, preserves its topological-layer profile and edge-to-node ratio, and then resamples node and edge costs using the same $(k,l,\mu,\alpha)$ model as the real graphs.
